\clearpage
\begin{proofbox}
	\textbf{\textcolor{CProof}{思路提示}}
	周期等式证明的核心是反复利用已知条件，通过迭代或代换，最终导出$f(x+T)=f(x)$的结构。推导中需注意定义域平移不变性（即若$x$在定义域内，则$x$与平移常数加减后的值仍在定义域内）。

	\medskip
	\begin{center}\begin{tikzpicture}[font=\small,>=stealth]
\draw[gray!40,rounded corners] (0,0) rectangle (14,7.3);
\node at (7,6.75) {例：$f(t)=\cos(\pi t)$，平移 $1$ 变号，平移 $2$ 恢复};
\begin{scope}[shift={(1.6,3.2)},x=2.2cm,y=1.6cm]
\draw[->,gray] (-.15,0)--(4.8,0) node[right]{$t$};
\draw[->,gray] (0,-1.25)--(0,1.3) node[left]{$y$};
\draw[blue,very thick,domain=0:4.4,samples=150] plot(\x,{cos(180*\x)});
\foreach \x/\y in {0/1,1/-1,2/1,3/-1,4/1}{\fill[orange!85!black] (\x,\y) circle(1.7pt);\ifnum\x>0\node[below=4pt] at (\x,0){$\x$};\fi}
\node[above right=5pt] at (0,1){$1$};\node[below=5pt] at (1,-1){$-1$};
\draw[<->,teal!80!black,thick] (.05,1.6)--(1.95,1.6) node[midway,above]{$T=2$};
\node[below left=4pt] at (0,0){$0$};
\end{scope}
\node at (7,.6) {$f(t+1)=-f(t)$，$f(t+2)=f(t)$；图中取 $a=1$};
\end{tikzpicture}
\end{center}
图以反号型为例展示两次平移后高度恢复；以下逐条用代入或换元核验。
\textbf{\textcolor{CProof}{正式推导}}
	\begin{description}[style=nextline, leftmargin=2.8em, labelsep=0.5em, itemsep=0.55em]
		\item[\textbf{步骤一（加常数直接型）：}]
			已知$f(x+a)=f(x)$恒成立。
			\[
				f(x+a)=f(x).
			\]
			\par\textbf{理由：} 若$a>0$，则直接取$T=a$满足$f(x+T)=f(x)$；若$a<0$，由$f(x+a)=f(x)$将$x$换为$x-a$得$f(x)=f(x-a)$，即$f(x+|a|)=f(x)$，故$|a|$是一个周期。
		\item[\textbf{步骤二（反号两倍型）：}]
			已知$f(x+a)=-f(x)$。计算$f(x+2a)$：
			\[
				\begin{aligned}
					f(x+2a) &= f((x+a)+a) \\ &= -f(x+a) \\ &= -(-f(x)) \\ &= f(x).
			\end{aligned}
			\]
			\par\textbf{理由：} 先令$x\mapsto x+a$代入条件得$f(x+2a)=-f(x+a)$，再利用原式替换$f(x+a)=-f(x)$，两次变号抵消，得到$f(x+2a)=f(x)$，故$2|a|$是周期。
		\item[\textbf{步骤三（倒数两倍型）：}]
			已知$f(x+a)=\pm\frac{1}{f(x)}$且$f(x)\neq0$。计算$f(x+2a)$：
			\[
				\begin{aligned}
					f(x+2a) &= f((x+a)+a) \\ &= \pm\frac{1}{f(x+a)} \\ &= \pm\frac{1}{\pm\frac{1}{f(x)}} \\ &= f(x).
			\end{aligned}
			\]
			\par\textbf{理由：} 连续两次应用倒数条件，无论正负，复合后均恢复原函数值，因此$2|a|$是周期。
		\item[\textbf{步骤四（差值平移型）：}]
			已知$f(x+a)=f(x+b)$（$a\neq b$）。令$t=x+b$，则$x=t-b$，代入得：
			\[
				f(t+a-b)=f(t).
			\]
			\par\textbf{理由：} 变量代换将两个不同平移的等式转化为标准形式$f(t+(a-b))=f(t)$，因此$|a-b|$是$f$的一个周期。
	\end{description}

	\medskip
	\textbf{\textcolor{CProof}{结论回扣}}
	以上四个步骤分别证明了输入中给出的四种常见周期判定结论。注意最终周期均取绝对值，且定义域平移不变性是所有结论成立的前提。在实际解题中，可根据条件形式直接套用相应周期公式。
\end{proofbox}

\clearpage
\begin{proofbox}
\textbf{新增证明一：把第一种复合分式完整展开}
\begin{center}\begin{tikzpicture}[font=\small,>=stealth]
\draw[gray!40,rounded corners] (0,0) rectangle (14,8.3);
\node at (7,7.7) {第一种分式：$2\mapsto3\mapsto2$，每次向右平移 $a=1$};
\begin{scope}[shift={(1.8,2.0)},x=2.3cm,y=1.25cm]
\draw[->,gray] (-.15,0)--(4.65,0) node[right]{$t$};
\draw[->,gray] (0,-.55)--(0,3.5) node[above]{$y$};
\draw[blue,very thick] (0,2)--(1,2);
\fill[blue] (0,2) circle(2pt);\draw[blue,fill=white] (1,2) circle(2pt);
\draw[orange!85!black,very thick] (1,3)--(2,3);
\fill[orange!85!black] (1,3) circle(2pt);\draw[orange!85!black,fill=white] (2,3) circle(2pt);
\draw[blue,very thick] (2,2)--(3,2);
\fill[blue] (2,2) circle(2pt);\draw[blue,fill=white] (3,2) circle(2pt);
\draw[orange!85!black,very thick] (3,3)--(4,3);
\fill[orange!85!black] (3,3) circle(2pt);\draw[orange!85!black,fill=white] (4,3) circle(2pt);
\foreach \x in {1,2,3,4}{\node[below=4pt] at (\x,0){$\x$};}
\node[left=6pt] at (0,2) {$2$};
\node[left=6pt] at (0,3) {$3$};
\fill[blue] (4,2) circle(2pt);
\node[below left=4pt] at (0,0){$0$};
\end{scope}
\node at (7,1.0) {取整段 $[2k,2k+1)$ 上为 $2$，$[2k+1,2k+2)$ 上为 $3$（$k\in\mathbb Z$）};
\node at (7,.45) {实心点取到、空心点不取到；这是满足条件的一个具体函数};
\end{tikzpicture}
\end{center}
令 $u=f(x)$，已知 $u\ne1$。先检查下一次的分母：
\[
 f(x+a)-1=\frac{u+1}{u-1}-1=\frac{2}{u-1}\ne0.
\]
再把条件中的 $x$ 换成 $x+a$，逐步代入：
\[
\begin{aligned}
 f(x+2a)
 &=\frac{f(x+a)+1}{f(x+a)-1}\\[5pt]
 &=\frac{\dfrac{u+1}{u-1}+1}{\dfrac{u+1}{u-1}-1}\\[5pt]
 &=\frac{\dfrac{u+1+u-1}{u-1}}{\dfrac{u+1-(u-1)}{u-1}}\\[5pt]
 &=\frac{\dfrac{2u}{u-1}}{\dfrac{2}{u-1}}=u=f(x).
\end{aligned}
\]
因此 $2a$ 是一个非零周期，取正值得 $\boxed{T=2|a|}$。负平移可由换元得到相反方向的正平移。
\end{proofbox}
\clearpage
\begin{proofbox}
\textbf{新增证明二：第二种分式注意括号与正负号}
\begin{center}\begin{tikzpicture}[font=\small,>=stealth]
\draw[gray!40,rounded corners] (0,0) rectangle (14,7.6);
\node at (7,7) {第二种分式：$2\mapsto-\tfrac13\mapsto2$，每次向右平移 $a=1$};
\begin{scope}[shift={(1.8,2.5)},x=2.3cm,y=1.25cm]
\draw[->,gray] (-.15,0)--(4.65,0) node[right]{$t$};
\draw[->,gray] (0,-.55)--(0,2.7) node[above]{$y$};
\draw[blue,very thick] (0,2)--(1,2);
\fill[blue] (0,2) circle(2pt);\draw[blue,fill=white] (1,2) circle(2pt);
\draw[orange!85!black,very thick] (1,-0.333333333333)--(2,-0.333333333333);
\fill[orange!85!black] (1,-0.333333333333) circle(2pt);\draw[orange!85!black,fill=white] (2,-0.333333333333) circle(2pt);
\draw[blue,very thick] (2,2)--(3,2);
\fill[blue] (2,2) circle(2pt);\draw[blue,fill=white] (3,2) circle(2pt);
\draw[orange!85!black,very thick] (3,-0.333333333333)--(4,-0.333333333333);
\fill[orange!85!black] (3,-0.333333333333) circle(2pt);\draw[orange!85!black,fill=white] (4,-0.333333333333) circle(2pt);
\foreach \x in {1,2,3,4}{\node[above=5pt] at (\x,0){$\x$};}
\node[left=6pt] at (0,2) {$2$};
\node[left=6pt] at (0,-0.3333333333333333) {$-\tfrac13$};
\fill[blue] (4,2) circle(2pt);
\node[above right=4pt] at (0,0){$0$};
\end{scope}
\node at (7,1.0) {取整段 $[2k,2k+1)$ 上为 $2$，$[2k+1,2k+2)$ 上为 $-1/3$（$k\in\mathbb Z$）};
\node at (7,.45) {实心点取到、空心点不取到；这是满足条件的一个具体函数};
\end{tikzpicture}
\end{center}
令 $u=f(x)$，已知 $u\ne-1$。下一次的分母同样不会为零：
\[
 1+f(x+a)=1+\frac{1-u}{1+u}=\frac{2}{1+u}\ne0.
\]
连续使用两次条件：
\[
\begin{aligned}
 f(x+2a)
 &=\frac{1-f(x+a)}{1+f(x+a)}\\[5pt]
 &=\frac{1-\dfrac{1-u}{1+u}}{1+\dfrac{1-u}{1+u}}\\[5pt]
 &=\frac{\dfrac{1+u-(1-u)}{1+u}}{\dfrac{1+u+1-u}{1+u}}\\[5pt]
 &=\frac{\dfrac{2u}{1+u}}{\dfrac{2}{1+u}}=u=f(x).
\end{aligned}
\]
所以 $\boxed{T=2|a|}$ 是一个正周期。
\end{proofbox}
